\subsubsection{矩阵维度与模块分工}
计算对象写作 $Z_{i,j}=\sum_{k=0}^{63}X_{i,k}W_{k,j}+b_j$，其中 $0\le i<99$、$0\le j<256$。这里的 99 是输入序列位置数，64 是点积长度，256 是输出通道数。一个任务覆盖整张输出矩阵，不能把一个列 tile 或一个行组算作独立完成的任务。矩阵乘法后使用该任务的移位参数转换至 signed12，完整保存所有输出位置。

输入矩阵共有 6336 个元素，权重共有 16384 个元素，偏置按 256 列广播。输入和权重的重排服务于阵列读取，不改变 $k$ 维度上的配对关系。输入第 $i$ 行与权重第 $j$ 列始终产生输出 $(i,j)$；后处理即使延后执行，也必须携带原 tile 和行组标识。KWT 的模型背景见文献\cite{kwt}，本实现的具体维度、位宽和任务边界以本工程冻结配置为准。

\begin{figure}[H]\centering
\begin{tikzpicture}[x=1cm,y=1cm]
\draw[fill=pale,draw=ink] (0,0) rectangle (2.5,3.2);
\draw[fill=accent!22,draw=ink] (3.8,1.5) rectangle (8.3,3.2);
\draw[fill=pale,draw=ink] (9.6,0) rectangle (14.1,3.2);
\node at (1.25,1.7) {$X$};\node at (6.05,2.35) {$W$};\node at (11.85,1.7) {$Z$};
\node at (3.15,2.3) {$\times$};\node at (8.95,2.3) {$\longrightarrow$};
\node[note] at (1.25,-.45) {$99\times64$};\node[note] at (6.05,1.05) {$64\times256$};
\node[note] at (11.85,-.45) {$99\times256$};
\node[note] at (6.05,.25) {偏置按列广播\\量化在完整点积之后};
\draw[accent,very thick] (9.6,2.75) rectangle (10.5,3.2);
\node[note] at (11.8,3.65) {着色框示意一次 $4\times16$ 输出块};
\end{tikzpicture}
\caption{矩阵维度和分块计算对象。矩形尺寸只表示数据关系，不按元素数量等比例绘制。}
\end{figure}

\begin{table}[H]\centering\small
\caption{计算链路中的模块职责及交接条件}
\begin{tabularx}{\linewidth}{p{34mm}YY}\toprule
模块 & 负责的数据或控制 & 对下一级的交接条件\\\midrule
配置与任务控制 & 基址、字数、epoch、状态 & 合法启动且输入边界已提交\\
DMA 与输入解码 & 共享 SRAM 字流、元数据 & 响应归属正确且解码数据有效\\
计算核心 & 配对读取、主项及行列修正 & 完整 32 对已累加，结果标识有效\\
共享后处理 & 偏置、移位、RNE、饱和 & 4 个结果及对应位置同时有效\\
固定格式写出 & 12 位字段到 32 位总线字 & 字缓冲有空间且输出索引合法\\
完成控制 & 头部、数据及最终标志 & 最后一笔写响应已返回\\\bottomrule
\end{tabularx}
\end{table}

上述分工使数值正确性与传输正确性可以分别检查。数值通路只消费已接受的输入；传输通路只对已接受的命令维护在途计数。局部模块的 valid 不等价于整机完成，输出写命令的发出也不等价于共享 SRAM 已经确认完成。
